% Part: applied-modal-logics
% Chapter: epistemic-logic
% Section: properties-accessibility

\documentclass[../../../include/open-logic-section]{subfiles}

\begin{document}

\olfileid{aml}{el}{acc}

\olsection{સુલભતા સંબંધો અને જ્ઞાનસંબંધી સિદ્ધાંતો}

સામાન્ય મોડલ તર્કપદ્ધતિઓમાં ફ્રેમ-અનુરૂપતા વિશે આપણે જે જાણીએ
છીએ તેને આધારે, જ્ઞાનસંબંધી અર્થઘટન આપતાં લાક્ષણિક
!!{formula}s કેવાં બને છે તે જોવું યોગ્ય છે. જ્ઞાનસંબંધી
તર્કપદ્ધતિઓનું અર્થઘટન સામાન્ય રીતે S5માં થાય છે એમ આપણે કહી
ચૂક્યા છીએ. હવે જુદા જ્ઞાનસંબંધી સિદ્ધાંતો કેવી રીતે દર્શાવાય
છે અને તેઓ જુદી ફ્રેમ-શરતોને કેવી રીતે અનુરૂપ છે તે જોઈએ.

સામાન્ય મોડલ તર્કમાંથી યાદ કરો કે જુદાં મોડલ !!{formula}s
સુલભતા સંબંધોના જુદા ગુણધર્મોને લક્ષણિત કરતાં હતાં. નીચેના
કોષ્ટકમાં ચોક્કસ જ્ઞાનસંબંધી સિદ્ધાંતોને અનુરૂપ એવા થોડા
ગુણધર્મો પસંદ કર્યા છે.

\begin{table}[t]
    \begin{tabular}{| p{.48\textwidth} || p{.48\textwidth} |}
      \hline
      {\emph{જો $R$ \dots હોય}} & {\emph{તો~$\mModel{M}$માં \dots સત્ય છે:}} \\
      \hline \hline
        
      & $\Knows (p \lif q) \lif (\Knows p \lif \Knows q)$ 
      \hfill \newline{(સંવરણ)} \\
      \hline
      \emph{સ્વવાચક}: $\forall w Rww$  
      & $\Knows p \lif p$ \hfill \newline{(સત્યાનુરૂપતા)} \\
      \hline
      \emph{પરંપરિત}: \newline
      $\forall u \forall v \forall w ((Ruv \land Rvw) \lif Ruw)$ & 
      $\Knows p \lif \Knows \Knows p$ \hfill 
           \newline{(હકારાત્મક આત્મનિરીક્ષણ)} \\
      \hline 
      \emph{યુક્લિડિયન}: \newline
      $\forall w \forall u \forall v ((Rwu \land Rwv) \lif Ruv)$ &  
      $\lnot \Knows p \lif \Knows \neg \Knows p$ \hfill 
        \newline{(નકારાત્મક આત્મનિરીક્ષણ)} \\
      \hline
    \end{tabular}
    \caption{ચાર જ્ઞાનસંબંધી સિદ્ધાંતો.}
    \ollabel{tab:four}
  \end{table} 

$T$ સ્વયંસિદ્ધાંતને અનુરૂપ સત્યાનુરૂપતાને આ સિદ્ધાંતોમાં ઘણી વાર
સૌથી ઓછો વિવાદાસ્પદ ગણવામાં આવે છે, કારણ કે તે કહે છે કે
!!a{formula} જાણીતું હોય તો તે સાચું હોવું જ જોઈએ. સંવરણ તથા
હકારાત્મક અને નકારાત્મક આત્મનિરીક્ષણ ઘણાં વધુ વિવાદાસ્પદ છે.

$K$ સ્વયંસિદ્ધાંતને અનુરૂપ સંવરણનો અર્થ એ છે કે કર્તાનું જ્ઞાન
પ્રેરણ હેઠળ બંધ છે. અમુક કિસ્સાઓમાં આ વિશ્વસનીય લાગે. દાખલા
તરીકે, મને ખબર હોઈ શકે કે હું વિક્ટોરિયામાં હોઉં, તો વૅન્કુવર
ટાપુ પર હોઉં. વિચિત્ર સંશયવાદી પરિસ્થિતિઓને બાજુએ રાખીએ, તો
મને ખબર છે કે હું વિક્ટોરિયામાં છું; તેથી મને વૅન્કુવર ટાપુ
પર હોવાની પણ ખબર છે એવું સૂચવાય છે. આ કિસ્સામાં મારા જ્ઞાનનું
તાર્કિક સંવરણ સહજ લાગે. બીજી તરફ, આપણે આપણા જ્ઞાનમાંથી નીકળતાં
બધાં પરિણામો વિશે હંમેશા વિચારતા નથી; તેથી બીજા કિસ્સાઓમાં
આ સિદ્ધાંત ઓછો સહજ લાગી શકે છે.

હકારાત્મક આત્મનિરીક્ષણને ક્યારેક KK-સિદ્ધાંત કહે છે: હું કંઈક
જાણતો હોઉં, તો મને ખબર હોય કે હું તે જાણું છું. આ 4
સ્વયંસિદ્ધાંતનો જ્ઞાનસંબંધી સમકક્ષ છે. એ જ રીતે નકારાત્મક
આત્મનિરીક્ષણ કહે છે કે હું કંઈક \emph{જાણતો ન હોઉં}, તો મને
ખબર હોય કે હું તે જાણતો નથી; આ 5 સ્વયંસિદ્ધાંતનો સમકક્ષ છે.
\sourcecorrection{OLMD-078}{મૂળમાં બંને સિદ્ધાંતોના વિરોધી
દૃષ્ટાંત માટે વ્યક્તિને હકીકતમાં માહિતી ખબર હોય એવું કહે છે.
નકારાત્મક આત્મનિરીક્ષણને ખંડિત કરવા માહિતી ખરેખર અજાણી હોવી
જરૂરી છે; તેથી નીચે બંને સ્થિતિઓ અલગ કરી છે.}
બંને માટે સામાન્ય વિરોધી દૃષ્ટાંતો મળી શકે: મને કોઈ બાબત ખબર
છે કે નહીં તેની ખાતરી ન હોય. મને તે ખરેખર ખબર હોય તો હકારાત્મક
આત્મનિરીક્ષણ નિષ્ફળ થઈ શકે; અને ખરેખર ખબર ન હોય તો નકારાત્મક
આત્મનિરીક્ષણ નિષ્ફળ થઈ શકે.


\end{document}
